\subsection{Necesidad de Seleccionar un Simulador Base}

La propuesta desarrollada en este trabajo exigió intervenir la representación del vector
de estado para integrar una estrategia de compresión sin pérdida de precisión y
evaluar su efecto sobre memoria, tiempo y exactitud numérica. En consecuencia, la
selección del simulador base constituyó una decisión metodológica central, ya que
condicionó la viabilidad de la integración, el alcance de la instrumentación
experimental y el esfuerzo técnico requerido para desarrollar la solución.

En los simuladores cuánticos de estado completo tipo Schrödinger, el vector de estado
es la estructura principal del cálculo y, a la vez, la fuente dominante del consumo de
memoria. Por ello, el simulador seleccionado debía permitir una intervención
suficientemente directa sobre esa representación, sin trasladar el problema hacia capas
internas difíciles de modificar o poco accesibles para observación experimental. Bajo
estas condiciones, se consideraron cinco alternativas representativas: QuEST,
Intel-QS, qsim, Quantum++ y TMFQSfullstate. Todas ellas pertenecen al paradigma de
simulación de estado completo, pero difieren en el grado de apertura de su
arquitectura, en la organización de su código y en el contexto de escalabilidad para el
cual fueron concebidas.